%% kapa: exam=1 round=35 year=2024 period=1 form=A subject=real_estate subject_ko=부동산학원론 seq=19 number=99
%% source: Q-Net 감정평가사 기출문제 (cst00309) artlSeq=5213229 file=raw/1008/5213229/1교시 문제지 원본.pdf
%% conversion: kordoc 4.24.51 -> tools/gen1.py (auto); review=proofread-sonnet; flags=none
\begin{question}{99}
토지의 특성과 감정평가에 관한 내용이다. ( )에 들어갈 것으로 옳은 것은?
\begin{framedtext}
○ ( ㄱ )은 장래편익의 현재가치로 평가하게 한다. \par
○ ( ㄴ )은 원가방식의 평가를 어렵게 한다. \par
○ ( ㄷ )은 개별요인의 분석과 사정보정을 필요하게 한다. \par
\end{framedtext}
\begin{choices}
  \item ㄱ: 영속성, ㄴ: 부증성, ㄷ: 개별성
  \item ㄱ: 개별성, ㄴ: 영속성, ㄷ: 부동성
  \item ㄱ: 영속성, ㄴ: 개별성, ㄷ: 부증성
  \item ㄱ: 부증성, ㄴ: 영속성, ㄷ: 개별성
  \item ㄱ: 영속성, ㄴ: 개별성, ㄷ: 부동성
\end{choices}
\end{question}
\answer{1}
